\documentclass[10pt,a4paper]{amsart}
\usepackage{amsmath,amssymb,amsthm,graphicx}
\usepackage{verbatim}
\newenvironment{reference}{\comment}{\endcomment}
%\newenvironment{reference}{}{}
\newenvironment{slogan}{\comment}{\endcomment}
\newenvironment{history}{\comment}{\endcomment}

% For commutative diagrams we use Xy-pic
\usepackage[all]{xy}

% We use 2cell for 2-commutative diagrams.
\xyoption{2cell}
% Theorem environments.
%
\theoremstyle{plain}
\newtheorem{theorem}[subsection]{Theorem}
\newtheorem{proposition}[subsection]{Proposition}
\newtheorem{lemma}[subsection]{Lemma}

\theoremstyle{definition}
\newtheorem{definition}[subsection]{Definition}
\newtheorem{example}[subsection]{Example}
\newtheorem{exercise}[subsection]{Exercise}
\newtheorem{situation}[subsection]{Situation}

\theoremstyle{remark}
\newtheorem{remark}[subsection]{Remark}
\newtheorem{remarks}[subsection]{Remarks}

\numberwithin{equation}{subsection}

% Macros
%
\def\lim{\mathop{\mathrm{lim}}\nolimits}
\def\colim{\mathop{\mathrm{colim}}\nolimits}
\def\Spec{\mathop{\mathrm{Spec}}}
\def\Hom{\mathop{\mathrm{Hom}}\nolimits}
\def\Ext{\mathop{\mathrm{Ext}}\nolimits}
\def\SheafHom{\mathop{\mathcal{H}\!\mathit{om}}\nolimits}
\def\SheafExt{\mathop{\mathcal{E}\!\mathit{xt}}\nolimits}
\def\Sch{\mathit{Sch}}
\def\Mor{\mathop{\mathrm{Mor}}\nolimits}
\def\Ob{\mathop{\mathrm{Ob}}\nolimits}
\def\Sh{\mathop{\mathit{Sh}}\nolimits}
\def\NL{\mathop{N\!L}\nolimits}
\def\CH{\mathop{\mathrm{CH}}\nolimits}
\def\proetale{{pro\text{-}\acute{e}tale}}
\def\etale{{\acute{e}tale}}
\def\QCoh{\mathit{QCoh}}
\def\Ker{\mathop{\mathrm{Ker}}}
\def\Im{\mathop{\mathrm{Im}}}
\def\Coker{\mathop{\mathrm{Coker}}}
\def\Coim{\mathop{\mathrm{Coim}}}

% Boxtimes
%
\DeclareMathSymbol{\boxtimes}{\mathbin}{AMSa}{"02}

%
% Macros for moduli stacks/spaces
%
\def\QCohstack{\mathcal{QC}\!\mathit{oh}}
\def\Cohstack{\mathcal{C}\!\mathit{oh}}
\def\Spacesstack{\mathcal{S}\!\mathit{paces}}
\def\Quotfunctor{\mathrm{Quot}}
\def\Hilbfunctor{\mathrm{Hilb}}
\def\Curvesstack{\mathcal{C}\!\mathit{urves}}
\def\Polarizedstack{\mathcal{P}\!\mathit{olarized}}
\def\Complexesstack{\mathcal{C}\!\mathit{omplexes}}
% \Pic is the operator that assigns to X its picard group, usage \Pic(X)
% \Picardstack_{X/B} denotes the Picard stack of X over B
% \Picardfunctor_{X/B} denotes the Picard functor of X over B
\def\Pic{\mathop{\mathrm{Pic}}\nolimits}
\def\Picardstack{\mathcal{P}\!\mathit{ic}}
\def\Picardfunctor{\mathrm{Pic}}
\def\Deformationcategory{\mathcal{D}\!\mathit{ef}}
\usepackage{cleveref}
\crefname{originallemma}{original lemma}{original lemmas}
\Crefname{originallemma}{Original lemma}{Original lemmas}
\makeatletter
\newlabel{categories-lemma-limits-products-equalizers}{{002N}{source}{}{sourcepointer-1}{}}
\newlabel{categories-lemma-limits-products-equalizers@cref}{{[originallemma][0][]002N}{[1][0][]source}{}{}{}}
\makeatother
\title{The category of differential graded modules is abelian and complete and cocomplete}
\author{The Stacks Project Authors}
\date{stacks-project:dga.tex:local-snapshot-20261012; GNU FDL 1.2 or later}
\begin{document}
\maketitle
\noindent\textit{Editorial note.} All delivered mathematical text is the author's own, copied from the pinned Stacks Project chapter file sources/stacks-dga/pinned-original-source.tex; no statement or proof step was rewritten, reordered or removed.\par\medskip
\noindent\textit{Editorial note.} The counted proof is the author's complete argument: the abelian and colimit claims proceed by creation through the forgetful functor (the source states the commutation of kernels, cokernels, sums and colimits with the forgetful functor); the product computation (prod N\_i)\textasciicircum{}n = prod N\_i\textasciicircum{}n with total object the direct sum of degrees is the substantive point and is displayed in the span; equalizers are covered since kernels are equalizers; the limits criterion is standard external tooling (stub).\par\medskip
\noindent\textit{Editorial note.} The label lemma-dgm-abelian names the counted statement, and the proof environment immediately following it in the source is delivered verbatim as its proof.\par\medskip
\noindent\textit{Editorial note.} The counted statement is carried by the author's own lemma environment in the chapter source.\par\medskip
\noindent\textit{Allowed standard prerequisites (editorial).}\begin{itemize}
\item The category of graded A-modules. Standard background: the category of (graded) A-modules is abelian with arbitrary limits and colimits. The creation argument transfers the abelian structure and the (co)limits from the category of A-modules through the forgetful functor.
\item Limits criterion. Categories, Lemma categories-lemma-limits-products-equalizers (standard tooling). Final step (source lines 435-437): a category with products and equalizers has all limits.
\end{itemize}\medskip
\ifdefined\hypertarget\hypertarget{sourcepointer-1}{}\fi
\noindent\textit{Source reference.} Original lemma 002N: sources/stacks-chow/tags.txt, source lines 104--104. Prerequisite pointer invoked at proof line(s) 436; it names a standard fact of a sibling chapter that the counted argument uses, and every step of the counted proof built on it is retained verbatim in the delivered body.\par\medskip
\begin{lemma}
\label{lemma-dgm-abelian}
Let $(A, d)$ be a differential graded algebra. The category
$\text{Mod}_{(A, \text{d})}$ is abelian and has arbitrary limits and colimits.
\end{lemma}

\begin{proof}
Kernels and cokernels commute with taking underlying $A$-modules.
Similarly for direct sums and colimits. In other words, these operations
in $\text{Mod}_{(A, \text{d})}$ commute with the forgetful functor to the
category of $A$-modules. This is not the case for products and limits.
Namely, if $N_i$, $i \in I$ is a family of
differential graded $A$-modules, then the product $\prod N_i$ in
$\text{Mod}_{(A, \text{d})}$ is given by setting $(\prod N_i)^n = \prod N_i^n$
and $\prod N_i = \bigoplus_n (\prod N_i)^n$. Thus we see that the product
does commute with the forgetful functor to the category of graded $A$-modules.
A category with products and equalizers has limits, see
Categories, Lemma \ref{categories-lemma-limits-products-equalizers}.
\end{proof}

\end{document}
